\documentclass[12pt]{ximera}
\title{Homework 7: Lone-Chooser, Sealed Bids, and Markers}
\begin{document}
\begin{abstract}
Sections 3.4--3.6: designing a lone-chooser split to meet a scheduling constraint, the
method of sealed bids for household goods and for chores, and the method of markers.
\end{abstract}
\maketitle

\section*{Sections 3.4--3.6: Lone-Chooser, Sealed Bids, and Markers}

\begin{problem}
Three research groups (A, B, and C) share a lab and need to divide $24$ hours of lab time
per day using the lone-chooser method. Group A doesn't care which hours it gets as long
as it has at least $8$ (a fair share), so it agrees to be the first divider.

Group A splits the day into $\set{\text{midnight--noon}}$ and
$\set{\text{noon--midnight}}$, and group B, preferring afternoons and evenings, takes
noon--midnight. For the next round, group A splits its half into midnight--4am, 4am--8am,
and 8am--noon. Group B will split its half into three parts it considers equal (not
necessarily consecutive hours), and then group C chooses one piece from A's split and one
from B's.

\begin{prompt}
Group C wants lab access for at least one hour at a time, at precise $8$-hour intervals
(such as 1--2am, 9--10am, and 5--6pm). Describe one way group B could split its portion
so that the final division satisfies all three groups.

\begin{freeResponse}
\end{freeResponse}

\begin{hint}
C gets only one of A's $4$-hour blocks, and each of those contains at most one time from
an $8$-hour pattern. So two of C's three slots must come from a \emph{single} piece of
B's split. Which patterns put two slots between noon and midnight?
\end{hint}

\begin{feedback}
C can take only one of A's blocks, so C's piece from B must contain \emph{two} of its
three hourly slots --- $8$ hours apart, both between noon and midnight. The pattern
5am, 1pm, 9pm works: 5--6am lies in A's 4am--8am block, and 1--2pm and 9--10pm both lie in
B's half.

So B could split its twelve hours into three $4$-hour parts, for example
\begin{itemize}
\item Part 1: 12--2pm and 9--11pm;
\item Part 2: 2--6pm;
\item Part 3: 6--9pm and 11pm--midnight.
\end{itemize}
C then takes Part 1 from B and 4am--8am from A, getting 5--6am, 1--2pm, and 9--10pm,
exactly $8$ hours apart. A keeps midnight--4am and 8am--noon, $8$ hours, which is all it
wanted. B keeps Parts 2 and 3, two of the three parts it considered equal, so at least
$\frac23$ of the half it valued at a fair share or more. Everyone is satisfied.

Many other splits work, as long as one of B's parts contains both 1--2pm and 9--10pm (or
a similar pair $8$ hours apart) and C can pair it with a matching block of A's.
\end{feedback}
\end{prompt}
\end{problem}

\begin{problem}
You and a few roommates have shared an apartment for the last two years, but now it's
time to move out, and you've agreed to divvy your shared belongings using the method of
sealed bids. The values you each assign are:

\begin{center}
\begin{tabular}{|r|c|c|c|c|}\hline
 & \textbf{You} & \textbf{Sam} & \textbf{Riley} & \textbf{Jordan} \\\hline
Sofa & \$205 & \$260 & \$185 & \$215 \\
Smart TV & \$210 & \$150 & \$180 & \$175 \\
Microwave & \$55 & \$75 & \$65 & \$50 \\
Coffee Maker & \$60 & \$25 & \$90 & \$75 \\
Blender & \$60 & \$30 & \$75 & \$100 \\
Vacuum Cleaner & \$120 & \$150 & \$130 & \$105 \\
Dining Table & \$270 & \$230 & \$255 & \$300 \\\hline
\end{tabular}
\end{center}

\begin{prompt}
\textbf{(a)} Which roommate gets to keep each item? (Each item goes to the highest
bidder.)

Sofa: $\answer{Sam}$ \quad Smart TV: $\answer{You}$ \quad Microwave: $\answer{Sam}$

\smallskip
Coffee Maker: $\answer{Riley}$ \quad Blender: $\answer{Jordan}$ \quad
Vacuum Cleaner: $\answer{Sam}$ \quad Dining Table: $\answer{Jordan}$

\begin{feedback}
Reading down each row for the largest bid:
\[
\begin{array}{lll}
\text{Sofa} & \text{Sam} & \$260\\
\text{Smart TV} & \text{You} & \$210\\
\text{Microwave} & \text{Sam} & \$75\\
\text{Coffee Maker} & \text{Riley} & \$90\\
\text{Blender} & \text{Jordan} & \$100\\
\text{Vacuum Cleaner} & \text{Sam} & \$150\\
\text{Dining Table} & \text{Jordan} & \$300
\end{array}
\]
\end{feedback}
\end{prompt}

\begin{prompt}
\textbf{(b)} What is each roommate's ``fair share'' in this system?

You: $\$\answer{245}$ \quad Sam: $\$\answer{230}$ \quad Riley: $\$\answer{245}$ \quad
Jordan: $\$\answer{255}$

\begin{feedback}
Each roommate's total valuation, divided by four:
\[
\begin{array}{lll}
\text{You}: & 980 & \to 980/4 = \$245\\
\text{Sam}: & 920 & \to 920/4 = \$230\\
\text{Riley}: & 980 & \to 980/4 = \$245\\
\text{Jordan}: & 1020 & \to 1020/4 = \$255
\end{array}
\]
\end{feedback}
\end{prompt}

\begin{prompt}
\textbf{(c)} Compare each roommate's fair share with their perceived value of the items
they've won. What does each owe, or what are they owed, in the initial settlement?
Enter a positive number for money received and a negative number for money paid in.

You: $\answer{35}$ \quad Sam: $\answer{-255}$ \quad Riley: $\answer{155}$ \quad
Jordan: $\answer{-145}$

\begin{feedback}
\[
\begin{array}{lcccl}
 & \text{items won} & \text{fair share} & \text{difference} & \\
\text{You} & \$210 & \$245 & -\$35 & \text{receives } \$35\\
\text{Sam} & \$485 & \$230 & +\$255 & \text{pays } \$255\\
\text{Riley} & \$90 & \$245 & -\$155 & \text{receives } \$155\\
\text{Jordan} & \$400 & \$255 & +\$145 & \text{pays } \$145
\end{array}
\]
(Sam's $\$485$ is the Sofa, Microwave, and Vacuum: $260+75+150$. Jordan's $\$400$ is
the Blender and Dining Table: $100+300$.)
\end{feedback}
\end{prompt}

\begin{prompt}
\textbf{(d)} Once every roommate has items and/or compensation worth exactly their fair
share, how much money is left in the system --- that is, what is the surplus?

$\$\answer{210}$ \qquad and each roommate's quarter of it is $\$\answer{52.50}$

\begin{feedback}
Money in from Sam and Jordan is $255 + 145 = \$400$; money out to You and Riley is
$35 + 155 = \$190$. The surplus is
\[ \$400 - \$190 = \$210, \]
which splits four ways as $\$210/4 = \$52.50$ each.
\end{feedback}
\end{prompt}

\begin{prompt}
\textbf{(e)} After splitting the surplus equally, what is each roommate's cash position
in the final settlement? Enter a positive number for money received and a negative
number for money paid.

You: $\answer{87.50}$ \quad Sam: $\answer{-202.50}$ \quad Riley: $\answer{207.50}$ \quad
Jordan: $\answer{-92.50}$

\begin{feedback}
Add each roommate's \$52.50 share of the surplus to their initial settlement:
\[
\begin{array}{ll}
\text{You}: & +35 + 52.50 = \text{receives } \$87.50 \text{, plus the Smart TV}\\
\text{Sam}: & -255 + 52.50 = \text{pays } \$202.50 \text{, keeps Sofa, Microwave, Vacuum}\\
\text{Riley}: & +155 + 52.50 = \text{receives } \$207.50 \text{, plus the Coffee Maker}\\
\text{Jordan}: & -145 + 52.50 = \text{pays } \$92.50 \text{, keeps Blender and Dining Table}
\end{array}
\]
The cash balances: $-87.50 + 202.50 - 207.50 + 92.50 = 0$. \checkmark

Every roommate ends up with more than their fair share --- each is ahead by exactly
\$52.50 in their own valuation. That surplus exists because the roommates disagree
about what things are worth.
\end{feedback}
\end{prompt}
\end{problem}

\begin{problem}
Jake and their partner Jessie live together and are using the method of sealed bids to
split responsibilities for chores. They assign \emph{negative} values to each chore,
based on how much they would be willing to pay the other to do it instead. A
less-negative value means the task is less repulsive to that person.

\begin{center}
\begin{tabular}{|r|c|c|}\hline
 & \textbf{Jake} & \textbf{Jessie} \\\hline
Run the dishwasher & $-\$2$ & $-\$3$ \\
Do the laundry & $-\$3$ & $-\$5$ \\
Plan meals for the week & $-\$4$ & $-\$9$ \\
Shop for groceries & $-\$2$ & $-\$8$ \\
Take out the trash & $-\$3$ & $-\$2$ \\
Clean the kitchen & $-\$5$ & $-\$6$ \\
Clean the bathroom & $-\$8$ & $-\$4$ \\
Water the houseplants & $-\$7$ & $-\$1$ \\\hline
\end{tabular}
\end{center}

\begin{prompt}
\textbf{(a)} Based on the listed bids, who will be assigned each task?

Dishwasher: $\answer{Jake}$ \quad Laundry: $\answer{Jake}$ \quad
Meal planning: $\answer{Jake}$ \quad Groceries: $\answer{Jake}$

\smallskip
Trash: $\answer{Jessie}$ \quad Kitchen: $\answer{Jake}$ \quad
Bathroom: $\answer{Jessie}$ \quad Houseplants: $\answer{Jessie}$

\begin{hint}
Each item still goes to the \emph{highest} bidder. With negative numbers, the highest
bid is the least negative one --- the person who minds the chore least.
\end{hint}

\begin{feedback}
Each chore goes to whoever values it highest, which with negative numbers means
whoever minds it least:
\[
\begin{array}{lll}
\text{Dishwasher} & -2 > -3 & \text{Jake}\\
\text{Laundry} & -3 > -5 & \text{Jake}\\
\text{Meal planning} & -4 > -9 & \text{Jake}\\
\text{Groceries} & -2 > -8 & \text{Jake}\\
\text{Trash} & -2 > -3 & \text{Jessie}\\
\text{Kitchen} & -5 > -6 & \text{Jake}\\
\text{Bathroom} & -4 > -8 & \text{Jessie}\\
\text{Houseplants} & -1 > -7 & \text{Jessie}
\end{array}
\]
\end{feedback}
\end{prompt}

\begin{prompt}
\textbf{(b)} What are Jake's and Jessie's fair shares in this system? (These will be
negative.)

Jake: $\answer{-17}$ \qquad Jessie: $\answer{-19}$

\begin{feedback}
Jake's total valuation of all the chores is
\[ -2-3-4-2-3-5-8-7 = -\$34, \]
so his fair share is $-34/2 = -\$17$.

Jessie's total is
\[ -3-5-9-8-2-6-4-1 = -\$38, \]
so her fair share is $-38/2 = -\$19$.

A fair share here means each should end up doing chores they find no worse than that
amount of unpleasantness.
\end{feedback}
\end{prompt}

\begin{prompt}
\textbf{(c)} Compare part (b) with each partner's perceived value of the tasks they've
been assigned. What do they each owe, or what are they owed, in the initial settlement?

Jake's allotment value: $\answer{-16}$ \qquad Jessie's allotment value: $\answer{-7}$

\smallskip
Initial settlement --- Jake pays $\$\answer{1}$ and Jessie pays $\$\answer{12}$.

\begin{feedback}
Jake takes the dishwasher, laundry, meal planning, groceries, and kitchen:
\[ -2-3-4-2-5 = -\$16. \]
Jessie takes the trash, bathroom, and houseplants:
\[ -2-4-1 = -\$7. \]

Compare each with the fair share:
\begin{itemize}
\item Jake: $-16$ versus a fair share of $-17$. His chore load is $\$1$ \emph{less}
      unpleasant than fair, so he pays $\$1$ into the system.
\item Jessie: $-7$ versus a fair share of $-19$. Her load is $\$12$ less unpleasant
      than fair, so she pays $\$12$ in.
\end{itemize}
Both partners come out ahead of their fair share on chores alone, so both pay in ---
which is what creates the surplus.
\end{feedback}
\end{prompt}

\begin{prompt}
\textbf{(d)} What is the surplus after the initial settlement?

$\$\answer{13}$ \qquad and each partner's half is $\$\answer{6.50}$

\begin{feedback}
Jake pays \$1 and Jessie pays \$12, so the system holds
\[ \$1 + \$12 = \$13, \]
and nothing is paid out. Split evenly, that is $\$13/2 = \$6.50$ each.
\end{feedback}
\end{prompt}

\begin{prompt}
\textbf{(e)} After splitting the surplus equally, how much money is each partner making
or paying? Enter a positive number for money received, negative for money paid.

Jake: $\answer{5.50}$ \qquad Jessie: $\answer{-5.50}$

\begin{feedback}
Subtract each partner's \$6.50 share of the surplus from what they owed:
\begin{itemize}
\item Jake: paid \$1, receives \$6.50 back, so he nets $+\$5.50$ --- he \emph{receives}
      \$5.50.
\item Jessie: paid \$12, receives \$6.50 back, so she nets $-\$5.50$ --- she
      \emph{pays} \$5.50.
\end{itemize}

So Jessie pays Jake \$5.50 a week. The final arrangement:
\begin{itemize}
\item \textbf{Jake} runs the dishwasher, does laundry, plans meals, shops for
      groceries, and cleans the kitchen, and receives \$5.50.
\item \textbf{Jessie} takes out the trash, cleans the bathroom, waters the houseplants,
      and pays \$5.50.
\end{itemize}
Check: Jake's position is $-16 + 5.50 = -\$10.50$ against a fair share of $-\$17$, so he
is \$6.50 ahead. Jessie's is $-7 - 5.50 = -\$12.50$ against $-\$19$, also \$6.50 ahead.
\checkmark
\end{feedback}
\end{prompt}
\end{problem}

\begin{problem}
Six grandchildren are dividing up their grandma's Ty beanie baby collection using the
method of markers. They lay the $60$ beanie babies out in a randomly ordered array,
numbered $1$ through $60$, then each individually marks off six consecutive segments
that are approximately equally valuable to them.

\begin{center}
\begin{tabular}{r|c|c|c|c|c|c}
& 1st share & 2nd share & 3rd share & 4th share & 5th share & 6th share \\\hline
Allan  & 1--7  & 8--14  & 15--29 & 30--36 & 37--52 & 53--60 \\
Brenna & 1--10 & 11--20 & 21--36 & 37--43 & 44--48 & 49--60 \\
Caleb  & 1--5  & 6--20  & 21--31 & 32--40 & 41--51 & 52--60 \\
Devon  & 1--8  & 9--17  & 18--25 & 26--37 & 38--53 & 54--60 \\
Emma   & 1--9  & 10--18 & 19--32 & 33--46 & 47--50 & 51--60 \\
Frank  & 1--14 & 15--21 & 22--28 & 29--40 & 41--56 & 57--60
\end{tabular}
\end{center}

Each grandchild's five markers sit at the end of their first five shares.

\begin{prompt}
\textbf{(a)} Which of their shares is each grandchild awarded via this method?

Round 1 (first markers) --- awarded to $\answer{Caleb}$, who takes items 1--$\answer{5}$.

\smallskip
Round 2 (second markers) --- $\answer{Allan}$ takes items 6--$\answer{14}$.

\smallskip
Round 3 (third markers) --- $\answer{Devon}$ takes items 15--$\answer{25}$.

\smallskip
Round 4 (fourth markers) --- $\answer{Frank}$ takes items 26--$\answer{40}$.

\smallskip
Round 5 (fifth markers) --- $\answer{Brenna}$ takes items 41--$\answer{48}$.

\smallskip
The last remaining grandchild, $\answer{Emma}$, takes items $\answer{51}$--60.

\begin{hint}
In each round, look only at the relevant marker (first markers in round 1, second
markers in round 2, and so on) among the players who have not yet been served, and
find the leftmost one that lies past the items already handed out.
\end{hint}

\begin{feedback}
The markers sit at the ends of the first five shares:
\[
\begin{array}{ll}
\text{Allan}: & 7,\ 14,\ 29,\ 36,\ 52\\
\text{Brenna}: & 10,\ 20,\ 36,\ 43,\ 48\\
\text{Caleb}: & 5,\ 20,\ 31,\ 40,\ 51\\
\text{Devon}: & 8,\ 17,\ 25,\ 37,\ 53\\
\text{Emma}: & 9,\ 18,\ 32,\ 46,\ 50\\
\text{Frank}: & 14,\ 21,\ 28,\ 40,\ 56
\end{array}
\]

\textbf{Round 1:} first markers are $7, 10, 5, 8, 9, 14$. Caleb's is leftmost at $5$,
so Caleb takes items $1$--$5$ and leaves.

\textbf{Round 2:} among the rest, second markers are Allan $14$, Brenna $20$,
Devon $17$, Emma $18$, Frank $21$. Allan's is leftmost, so Allan takes $6$--$14$.

\textbf{Round 3:} third markers are Brenna $36$, Devon $25$, Emma $32$, Frank $28$.
Devon takes $15$--$25$.

\textbf{Round 4:} fourth markers are Brenna $43$, Emma $46$, Frank $40$. Frank takes
$26$--$40$.

\textbf{Round 5:} fifth markers are Brenna $48$ and Emma $50$. Brenna takes $41$--$48$.

\textbf{Emma}, the last one left, takes everything past her final marker at $50$:
items $51$--$60$.
\end{feedback}
\end{prompt}

\begin{prompt}
\textbf{(b)} Which beanie babies remain after every grandchild has been awarded exactly
their fair share?

Items $\answer{49}$ and $\answer{50}$ are left over.

\begin{feedback}
The awarded segments are $1$--$5$, $6$--$14$, $15$--$25$, $26$--$40$, $41$--$48$, and
$51$--$60$. The only numbers not covered are $49$ and $50$.
\end{feedback}
\end{prompt}

\begin{prompt}
\textbf{(c)} Describe one way that the rest of the beanie babies can be fairly
distributed.

\begin{freeResponse}
\end{freeResponse}

\begin{feedback}
Every grandchild already holds a share they consider fair, so anything done with the
two leftovers only improves someone's position --- there is no way to make the division
unfair at this stage. Reasonable options include:

\begin{itemize}
\item \textbf{Run the method again} on the leftover items. With only two items and six
      players this is awkward, but it is the textbook answer when many items remain.
\item \textbf{Draw lots}, giving each leftover to a randomly chosen grandchild, with no
      one receiving two.
\item \textbf{Auction them} among the six, splitting the proceeds equally --- effectively
      the method of sealed bids applied to the remainder.
\item \textbf{Let the grandchildren take turns choosing}, in an order decided randomly.
\end{itemize}

Any of these is acceptable. What matters is that the leftovers are distributed by some
process the players agree to in advance, since the fairness guarantee of the method of
markers has already been satisfied.
\end{feedback}
\end{prompt}
\end{problem}

\end{document}
